% =====================================================================
%  BAB 2 — METODE
% =====================================================================

\chapter{Metode}

\section{Dataset}
\label{sec:2.1}

Karakteristik dataset dirangkum pada Tabel~\ref{tab:1.data}.

\begin{thesistable}
\begin{longtable}{@{}l r r p{4cm}@{}}
  \caption{Karakteristik dataset}
  \label{tab:1.data}\\
  \toprule
  \textbf{Subset} & \textbf{Jumlah} & \textbf{Proporsi (\%)} & \textbf{Keterangan} \\
  \midrule
  \endfirsthead
  \multicolumn{4}{@{}l}{\footnotesize\itshape Lanjutan Tabel~\thetable}\\
  \toprule
  \textbf{Subset} & \textbf{Jumlah} & \textbf{Proporsi (\%)} & \textbf{Keterangan} \\
  \midrule
  \endhead
  \bottomrule
  \endfoot
  \bottomrule
  \endlastfoot
  {[Kelas A]} & [X.XXX] & [XX,X] & [Deskripsi] \\
  {[Kelas B]} & [X.XXX] & [XX,X] & [Deskripsi] \\
  \midrule
  \textbf{Total} & [X.XXX] & 100,0 & \\
\end{longtable}
\end{thesistable}

\section{Tabel Lebar}
\label{sec:2.2}

\begin{thesistable}
\begin{landscape}
\begin{longtable}{@{}l r@{}}
  \caption{Rekapitulasi}
  \label{tab:2.rekap}\\
  \toprule
  \textbf{Kategori} & \textbf{Nilai} \\
  \midrule
  \endfirsthead
  \multicolumn{2}{@{}l}{\footnotesize\itshape Lanjutan Tabel~\thetable}\\
  \toprule
  \textbf{Kategori} & \textbf{Nilai} \\
  \midrule
  \endhead
  \bottomrule
  \endfoot
  \bottomrule
  \endlastfoot
  A & 1,5 \\
  B & 2,5 \\
\end{longtable}
\end{landscape}
\end{thesistable}
